% Archived source excerpt from v76; preserved unchanged below.
% Not included in the v77 compiled SI; claims require the repairs recorded in the v77 implementation record.

\section{Framework extensions in full}\label{si:framework}

The report editions of this paper carried, in a separate markdown supplement, a set of framework extensions that the main article does not use in any proof: the full named record behind the typed framework, a governance extension in which instruments enter as constructors with declared support, an intergenerational and a composition extension, an application note on planetary boundaries and a register of declared conjectures. They are reproduced here in full, with the status each item carries (definition, identity, theorem, conditional theorem, conjecture or counterexample), so that the research-article edition loses none of the deposited material. Nothing in this section is load-bearing for the main-text results; the cross-references to the main text use the numbering of the present edition.

\subsection{Extended framework definitions}

The framework's full definitions are collected here. The named record $\mathfrak{S} = (T, Z, S_{st}, B_{out}, V, \Gamma, O, A, C, R, D, K, P)$ has thirteen fields, each of which a sustainability-assessment model must declare before any invariance, viability, or aggregation claim can be transferred between models. The four uncertainty levels and the four model maps complete the definitional layer.

\emph{Notation.} We write $\mathfrak{S} = (T, Z, S_{st}, B_{out}, V, \Gamma, O, A, C, R, D, K, P)$ for the named record; its slots are itemized below. A constructor -- a map on the system data representing a governance instrument -- is written $\mathsf{C}$, with declared support $\operatorname{supp}(\mathsf{C})$. Existential viability kernels are written $\mathrm{Viab}(V; U, F)$.

The accepted-state notation $\mathcal{V}[\cdot]$ follows the main text and is distinct from the policy class $P$. Generation-indexed constraint sets are written $V^{(k)}$.

\begin{itemize}[leftmargin=2em]
\item \textbf{Type system $T$.} Named types with units; moieties (conserved substances) are the only types admitting addition. Bridge types (services, thresholds, information states, institutional variables) are distinct.
\item \textbf{State space $Z$.} Typed product of stock, service, and information coordinates.
\item \textbf{Stock--flux structure $S_{st}$.} Per-moiety conservation laws; typed flux maps between stocks of the same moiety.
\item \textbf{Boundary interface $B_{out}$.} Declared external fluxes (harvest, emission, recharge) with owners.
\item \textbf{Constitutive laws $V$.} Growth, conversion, and transformation laws on typed stocks.
\item \textbf{Service--technology correspondence $\Gamma$.} Which services each architecture delivers and at what intensity.
\item \textbf{Observation operator $O$.} What each decision authority observes, with declared error structure.
\item \textbf{Assessment operator $A$.} The mapping from observations to verdicts; the subject of the main article.
\item \textbf{Command architecture $C$.} Who may select which actions, on what information, at what timing.
\item \textbf{Deployment/reset architecture $R$.} Endpoint resets and their costs (as in the witness datum).
\item \textbf{Disturbance class $D$.} Declared set of admissible disturbance trajectories.
\item \textbf{Safe-and-just set $K$.} The constraint set, with floors typed by provenance.
\item \textbf{Policy class $P$.} The set of causal feedback policies; distinct from the accepted-state notation $\mathcal{V}[\cdot]$.
\end{itemize}

\textbf{Uncertainty levels.} Four levels, each carrying a fixed quantifier discipline; no claim mixes levels without a declared bridge.

\begin{itemize}[leftmargin=2em]
\item (U1) Parameter uncertainty: fixed-but-unknown parameters in known structure.
\item (U2) Observation/assessment uncertainty: error between state and assessment inputs.
\item (U3) Process disturbance: the declared class $D$.
\item (U4) Structural model uncertainty: the model itself is one of several candidates.
\end{itemize}

\textbf{Model maps.} The four maps are not interchangeable; claims transfer only along declared maps.

\begin{itemize}[leftmargin=2em]
\item (M1) Specialisation: a model is an instance of a more general model class.
\item (M2) Exact projection: a semiconjugacy preserving the relevant invariants.
\item (M3) Approximation: a declared residual bound.
\item (M4) Singular reduction: a limit with declared validity region.
\end{itemize}

\subsubsection{Claim status}

Every statement carries a status: definition, identity, theorem, conditional theorem, conjecture, or counterexample. Two rules govern the vocabulary.

\begin{itemize}[leftmargin=2em]
\item \textbf{No promotion.} A conditional theorem is never stated as a theorem.
\item \textbf{No silent transfer.} A status established for one model class does not transfer to another without a declared map.
\end{itemize}

The status lines of this section apply the vocabulary item by item.

\subsubsection{Typed failure classes}

An assessment verdict must name the failure class it establishes. The taxonomy is definitional, not predictive: (1) material inconsistency; (2) physical infeasibility; (3) epistemic/common-prescription infeasibility; (4) authority infeasibility; (5) implementation infeasibility; (6) temporal infeasibility; (7) recovery failure; (8) architecture-transition failure; (9) model-credibility failure; (10) normative incompatibility. The discipline is to identify the earliest discharged obstruction rather than relabel an institutional failure as physical impossibility.

\subsubsection{The witness as a specialization}

The main article's witness datum (main-text Section~4) instantiates the record field by field; the mapping is recorded here so that the datum is a specialization of the framework rather than parallel text. On the witness: $T$ = the three moieties (reserve stock, protected-group service surplus, remediation-liability coverage) with their units; $Z$ = the phase space $(q, x, s_1, s_2)$; $S_{st}$ = the linear spend/growth flows of main-text Section~4; $B_{out}$ = none (closed interval); $V$ = the piecewise-linear tube laws; $\Gamma$ = the four named actions; $O$ = the exact state observation; $A$ = the typed assessment operators of main-text Section~2; $C$ = the single-command architecture (one action per interval); $R$ = the destination reset $s \mapsto s + e$; $D$ = $\{\alpha, \beta\}$ as specified in main-text Section~4; $K = S_0$; $P$ = the four-element menu.

The remaining record slots -- multi-command architectures, partially observed states, and policy classes larger than the menu -- are not used by the witness and are stated only as framework scope.

For composite-index construction, the extended definitions mean that any aggregation claim carries a tuple of declared slots and a declared uncertainty level; absent those declarations, the claim has no fixed meaning and cannot be checked against a typed assessment.

\subsection{Governance extension: constructors with declared support}

\noindent\textbf{Status.} Framework extension; the constructor algebra and the implementability ladder are stated at partial status. The obstruction theorem is a conditional result summarised from the companion obstruction calculus (Abaee, 2026a), with the proof conditions made explicit.

\subsubsection{Constructors}

Governance instruments enter the framework as \textbf{constructors}: maps on the system data representing a governance instrument. A constructor does not, in general, alter exactly one tuple component.

A harvest cap may alter the action correspondence, the dynamics, and the enforcement model; a tax may alter the effort law, the policy class, and the equilibrium; an observation reform may alter both the observation operator and the feedback law. Each primitive constructor $\mathsf{C}$ therefore carries a declared support

\[\operatorname{supp}(\mathsf{C}) \;\subseteq\; \{ T, Z, S_{st}, B_{out}, V, \Gamma, O, A, C, R, D, K, P \},\]

the set of tuple components it may alter. Any induced change outside the declared support must be represented explicitly or prohibited.

The declaration makes constructor composition auditable. For constructors $\mathsf{C}_1, \mathsf{C}_2$, one may prove commutation or noncommutation

\[\mathsf{C}_1 \circ \mathsf{C}_2 \;\neq\; \mathsf{C}_2 \circ \mathsf{C}_1\]

when, for example, an observation change alters the implementable policy class before versus after an action restriction.

The primitive constructors are: $\mathsf{Cap}(Q)$ (upper harvest bound), $\mathsf{Floor}(H_{\min})$ (output floor), $\mathsf{Tax}(\tau)$ (price entry in the effort law), $\mathsf{Excl}$ (excluding a competing predator or agent), $\mathsf{Leak}(h)$ (unreported/illegal harvest), $\mathsf{Obs}(I, \Psi)$ (replacing the observation map and feedback law), $\mathsf{Rest}(\cdot)$ (any other restriction of the control correspondence). Each must declare its support.

\subsubsection{Action-set monotonicity and the implementability ladder}

Restriction of the action correspondence obeys one-sided monotonicity: $U_1(x) \subseteq U_2(x)$ for every $x$ implies $\mathrm{Viab}(V; U_1, F) \subseteq \mathrm{Viab}(V; U_2, F)$ for existential viability under the same dynamics and information structure.

The family of correspondences with nonempty existential viability kernel is upward closed and closed under unions; intersections are not generally preserved.

This monotonicity concerns the \textbf{action correspondence}. It is separate from monotonicity of the \textbf{policy class}, which also depends on information, timing, authority, strategic equilibrium, observability, and implementation mechanisms.

The action correspondence is layered along an \textbf{implementability ladder} -- a chain of nested action correspondences ordered by what an agent can actually do at each layer of abstraction:

\[U_{\mathrm{impl}}(z) \;\subseteq\; U_{\mathrm{inst}}(z) \;\subseteq\; U_{\mathrm{tech}}(z) \;\subseteq\; U_{\mathrm{theor}}(z),\]

with the parallel policy-class ladder $P_{\mathrm{impl}} \subseteq P_{\mathrm{inst}} \subseteq P_{\mathrm{tech}} \subseteq P_{\mathrm{theor}}$.

Within-architecture viability quantifies over the implementable class. A viability result over a technological or theoretical class does not transfer downward, because restriction of the correspondence can empty the kernel even where enlargement preserves it.

\subsubsection{Management vocabularies as representations}

Some instruments admit a canonical representation by primitive constructors; the representation is valid only after proving the induced component changes. A total allowable catch is representable as $\mathsf{Cap}$; a harvest control rule as $\mathsf{Obs}$; a landing subsidy as $\mathsf{Tax}(-\sigma)$; unreported harvest as $\mathsf{Leak}$; a closed season as a periodic $\mathsf{Cap}$. Each representation is valid conditional on its declared support being discharged.

The representation claim is a theorem about the instrument, not a rename. A tax is not merely a price entry if it changes incentives, equilibrium controls, compliance, enforcement, revenue and redistribution, or the disturbance law; a predator exclusion can alter trophic dynamics; an observation reform can change the admissible policy class.

\subsubsection{The commons obstruction (finite-time result with boundary-exit condition)}

\begin{theorem}[commons obstruction; finite-time result with boundary-exit condition]\label{thm:si:commons}
Suppose the following hypotheses hold.
\begin{itemize}[leftmargin=3em]
\item[(H1)] A scalar stock $S$ obeys the dynamics $\dot S = g(S) - H$.
\item[(H2)] Under the implemented Nash feedback there exist constants $a > 0$ and $\varepsilon > 0$ such that
$g(S) - H^{\mathrm{Nash}}(S) \le -\varepsilon$ for every $S \in [S_{\min}, S_{\min} + a]$.
\item[(H3)] Every relevant trajectory enters this strip.
\end{itemize}
Then every trajectory entering the strip reaches $S_{\min}$ within time at most $a/\varepsilon$. If, in addition, the same strict negative-drift condition persists at the boundary and solutions are forward complete, the trajectory subsequently exits below $S_{\min}$.
\end{theorem}

\noindent\noindent\textbf{Status.} Summarised from the companion obstruction calculus (Abaee, 2026a), with the proof conditions made explicit. The uniform margin $\varepsilon > 0$ is load-bearing: strictly negative drift without a uniform margin does not by itself imply finite-time exit. The theorem requires a defined scalar stock dynamics and an explicit quantifier over Nash trajectories. Reaching the boundary is not identical to exiting below it; the second sentence requires the additional boundary condition.

\begin{remark}[monitoring is not automatic] Observation error $\|\hat x - x\| \le \delta$ does not by itself make a desired feedback law implementable. One needs a robust-feedback theorem certifying a control law whose robust admissibility holds for all states in the observation fibre. Otherwise monitoring may improve knowledge without producing a safe common prescription.\end{remark}

For governance design, the constructor extension means that instrument names must be discharged by support declarations before any invariance or viability claim can be transferred, and that finite-time obstruction results require an explicit uniform margin rather than a pointwise drift condition.

\subsection{Intergenerational extension}

\noindent\textbf{Status.} Framework extension. The nested-impossibility theorem is stated with its proof conditions made explicit. The ``every finite segment satisfiable'' phenomenon requires a separate example and is not implied by the theorem as stated.

\subsubsection{Generation structure}

A generation structure is a sequence $0 = t_0 < t_1 < \cdots \to \infty$ of generation boundaries with closed per-generation constraint sets $V^{(k)}$. Intergenerational viability requires $x(t) \in V^{(k)}$ for $t \in [t_k, t_{k+1})$.

\subsubsection{Stationary equivalence (with conditions)}

If $V^{(k)} = V$ for all $k$, intergenerational viability reduces to ordinary viability. The equivalence holds only under several restrictions: the generation intervals must cover the same infinite horizon; the policy class must be unchanged across generations; there must be no resets or inherited variables; the equality $V^{(k)} = V$ must be exact; and the intergenerational criterion must be simply $x(t) \in V^{(k)}$. If inheritance, reset, age structure, or generation-specific policy restrictions exist, stationary set equality alone may not reduce the model to ordinary viability.

\subsubsection{Nested-impossibility theorem}

\begin{theorem}[nested impossibility]\label{thm:si:nested}
Assume the following hypotheses.
\begin{itemize}[leftmargin=3em]
\item[(H1)] Trajectories are forward complete.
\item[(H2)] Every trajectory is confined to one common compact set $K$.
\item[(H3)] Each $V^{(k)}$ is closed.
\item[(H4)] The generation sets are nested decreasing with empty intersection, $\bigcap_k V^{(k)} = \varnothing$.
\end{itemize}
Then no intergenerationally viable path exists.
\end{theorem}
\begin{proof}
Suppose a viable path $x(t)$ exists with $x(t) \in V^{(k)}$ for all $t \ge t_k$. Since the trajectory is confined to the compact set $K$ and is forward complete, it has an $\omega$-limit point $x^\ast$. Because each $V^{(k)}$ is closed and the trajectory lies in $V^{(k)}$ for all sufficiently large $t$, $x^\ast\in V^{(k)}$ for every $k$. Hence $x^\ast\in\bigcap_k V^{(k)}$, contradicting emptiness.
\end{proof}

\textbf{What the theorem does not establish.} It proves no infinite path exists when the intersection is empty. It does not prove that every finite horizon is satisfiable. The phenomenon

\[\mathrm{Viab}_N \neq \varnothing \ \ \forall N, \qquad \mathrm{Viab}_\infty = \varnothing\]

requires a separate explicit example -- for instance $V^{(k)} = [0, 1/k]$ with dynamics and initial conditions making every finite prefix feasible but no point belonging to all sets. Until such an example is supplied, the claim is only the infinite-horizon impossibility.

For long-horizon sustainability assessment, the intergenerational extension means that intergenerational viability claims cannot be lifted from finite-horizon satisfiability alone; a separate argument is required whenever generations carry distinct constraint sets.

\subsection{Composition extension}

\noindent\textbf{Status.} Framework extension at conditional / incomplete status. The displayed formulas require equations and proofs that are not provided here; each item is marked.

\subsubsection{Effort-scale invariance (requires equations)}

In the registered delay-dynamics family of the companion study on institutional delay (Abaee, 2026e), the effort-scale transformation $E' = aE$, $E'_{\max} = aE_{\max}$, $q' = q/a$, $\delta_0' = a\delta_0$ is claimed to leave the $(N, Z)$ trajectories invariant.

\textbf{Status: conditional.} Invariance cannot be assessed without the registered equations: the transformed equations must be displayed and verified term by term, and the paper must distinguish parameter nonidentifiability, state-coordinate scaling, observational equivalence, and dynamical conjugacy, which are not identical. The architectural role is identifiability discipline: effort scales are not separately identifiable from the stock dynamics, so calibration claims must quotient by the transformation.

\subsubsection{Yield-gap soft-minimum (unsupported as written)}

At the Liebig limit, the yield gap is claimed to obey $\pi_j \le w_{\min}^{-1} e^{-\rho \Delta_y}$, with the coupled system decoupling with error $\|X - X^k\| \le C_T \varepsilon_c$, $\varepsilon_c = C e^{-\rho \Delta_y} + \varepsilon_{\mathrm{phys}}$.

\textbf{Status: conditional / unsupported as written.} This requires definitions of $w_{\min}$, $\pi_j$, $\Delta_y$, the yield functions, the normalization, the domain, and the assumptions under which the exponential estimate holds; the error estimate requires a stability or Lipschitz theorem, initial-condition matching, a finite time horizon, and a specified norm. It should be converted into a proposition with a complete assumption list and proof, or marked as a conditional source result.

\subsubsection{Coupling creates viability (not reproducible as written)}

An example with $g_i(s) = s(1-s)$ and coupling $d = 0.2$, referring to ``equilibrium-defined harvest floors,'' is claimed to show that coupling can create viability.

\textbf{Status: not reproducible as written.} To stand as a formal result it requires the coupled equations, the control/action set, the disturbance set, the safe set, the exact floor definition, the decoupled comparison system, the kernel calculation, and a proof or exact numerical certificate. Until supplied, this belongs in a motivating example, not in the formal framework. The two-sided lesson -- composition can create or destroy viability -- is retained as the architectural point, with unrestricted composition not licensed.

\subsubsection{Exergy, quality grades, and nonsmooth transformation feasibility}

\textbf{Status: declared research programme, not a theorem set.} Transformation feasibility under exergy and quality-grade constraints, where the feasible set is nonsmooth (grades induce kinks), is the declared frontier of the transformation operator. No status is asserted beyond programme.

For composite-index construction under composition, the composition extension means that composition can either create or destroy viability, and that any invariance claim across composed subsystems must be supported by equations and proofs rather than by assertion.

\subsection{Application note: planetary boundaries}

\noindent\textbf{Status.} Interpretive note; no theorem. The main article makes no claim about planetary boundaries.

The typed operator $E_{\mathrm{typ}}$ provides a formal idealization of the separate-boundary reading underlying the Planetary Boundaries framework (Rockstr\"om et al., 2009; Steffen et al., 2015), which is organized around nine boundary processes in the current standard formulation, with additional proposed indicators and extensions in the wider literature. The framework's operational logic is reasonably read as treating each boundary separately.

Planetary-boundary assessments also involve uncertainty, control variables, regional disaggregation, interactions, justice considerations, aggregation of risk, and boundary-transgression severity. Boundary processes do not act independently: Earth-system interactions amplify human impacts and shrink the safe operating space (Lade et al., 2020). The same separately-checked-floor logic underlies the safe-and-just-space (``doughnut'') literature, which evaluates social thresholds and biophysical ceilings independently, without cross-compensation, and reports that no country satisfies all floors within all ceilings (O'Neill et al., 2018; Fanning et al., 2022). We therefore state only that the typed operator provides a formal model of the separate-boundary reading, not that it captures the framework's full operational logic.

\textbf{What the main theorem does and does not contribute.} Main-text Theorem~2 shows that a particular compensatory assessment can accept a transition rejected by a typed assessment. It provides a formal model of \emph{one reason a noncompensatory stance may be adopted} -- namely, that per-weight compensation can license weight-dependent plans no single one of which respects all floors. It does not prove that planetary boundaries ought to be noncompensatory; that is a normative and scientific-design question. Nor does the witness's rescue mechanism generalize automatically: the witness has an explicit STAGED action funded by a scalar stock $x$, whereas many boundary violations have no single identifiable bridging resource. The rescue implication is conditional on the presence of a typed rescue action controlled by a resource margin.

For composite-index construction in the planetary-boundaries setting, this note means that the typed operator supplies a formal model of the separate-boundary reading only; whether the framework ought to be noncompensatory is a normative and scientific-design question that the main theorem does not settle.

\subsection{Declared conjectures and falsification designs}

\noindent\textbf{Status.} Conjecture and design status; none executed; none carries an empirical finding.

Nine conjectures:

\begin{itemize}[leftmargin=2em]
\item (C1) compositional sustainability -- viability is preserved under the composition of two viable subsystems exactly when their interaction map lies in a declared class;
\item (C2) transformability -- architecture transitions admit a resource-cost characterization analogous to the rescue threshold $\kappa^*$ of main-text Proposition~5;
\item (C3) capacity-leading failure -- the binding floor in an assessment sequence is the one with smallest slack at the earliest discharge;
\item (C4) bottleneck--robustness -- robustness margins concentrate in bottleneck subsystems;
\item (C5) boundary-expansion reversal -- expanding a constraint boundary can contract the viability kernel if the expansion admits a new attractor;
\item (C6) distributional dynamics -- the acceptance gap of main-text Theorem~2 is largest when floors are negatively correlated across the disturbance class;
\item (C7) correlated-disturbance amplification -- correlated disturbances across moieties amplify the gap;
\item (C8) maintenance suppression -- aggregate indices suppress maintenance investment signals relative to per-floor reporting;
\item (C9) efficiency--scale interaction -- the rescue threshold $\kappa^*$ is nonconvex in system scale.
\end{itemize}

Each conjecture is governed by preregistration restrictions. No conjecture is rescued by arbitrary post-hoc state augmentation. Each study preregisters system class, specification, candidate indicators, excluded variables, predicted direction, acceptable model revisions, and the observations that count against the conjecture. The unrestricted claim that every sustainability failure is representable at an ``adequate scale and resolution'' is excluded as too elastic to falsify. Each candidate leading indicator's predictive advantage over simpler outcome indicators is an empirical requirement, not a guarantee of definition.

For sustainability-assessment research design, the conjecture register means that conjectures must be preregistered with their falsification conditions before any empirical claim of leading-indicator advantage can be entertained.

